\section{A sharper oracle inequality for a Gaussian weight prior}
\label{sec:gaussian-oracle}

This section concerns a fully connected network with $H\geq1$ hidden
layers, all of width $N$, and a scalar linear readout.  Its total depth,
counting the readout, is $L=H+1$.  The activation in every hidden layer is
$\tanh$.  Writing $\bar x=(1,x)\in\mathbb R^{d+1}$, the network is
\[
 h^1(x)=\tanh(W^1\bar x),\qquad
 h^\ell(x)=\tanh(W^\ell h^{\ell-1}(x))\quad(2\leq\ell\leq H),
 \qquad z_w(x)=v^\top h^H(x).
\]
The activation is applied coordinatewise.  There are no internal biases;
the constant input coordinate supplies the first-layer biases.  Define
the weighted quadratic energy
\begin{equation}
 Q_N(w)=\frac dN\|W^1\|_F^2+
       \sum_{\ell=2}^H\|W^\ell\|_F^2+N\|v\|_2^2.
 \label{eq:gaussian-energy}
\end{equation}
Here and below $d\geq1$.  Inputs may have any distribution on
$\mathbb R^d$; the labels satisfy $|Y|\leq1$.  No bound on the input norm
is needed for the argument in this section.

Let $S=((X_i,Y_i))_{i=1}^n$ consist of independent samples from a law $D$.
For a predictor $g$, write
\[
 R(g)=\mathbb E_D(g(X)-Y)^2,\qquad
 \widehat R_S(g)=\frac1n\sum_{i=1}^n(g(X_i)-Y_i)^2,
 \qquad \operatorname{clip}(t)=\max\{-1,\min\{1,t\}\}.
\]
For positive $\sigma$ and $\lambda$, give the weights independent
centered Gaussian priors with variances
\begin{equation}
 \operatorname{Var}(W^1_{jk})=\frac{N\sigma^2}{d},\qquad
 \operatorname{Var}(W^\ell_{jk})=\sigma^2\ (2\leq\ell\leq H),\qquad
 \operatorname{Var}(v_j)=\frac{\sigma^2}{N}.
 \label{eq:gaussian-prior-variances}
\end{equation}
Use a Gaussian likelihood with mean $z_w(X_i)$ and variance
$n\lambda\sigma^2$ for each observation.  The posterior is therefore
\begin{equation}
 \Pi_{\sigma,N,S}(dw)=\frac1{Z_{\sigma,N,S}}
  \exp\left[-\frac{Q_N(w)+\lambda^{-1}\widehat R_S(z_w)}{2\sigma^2}\right]dw.
 \label{eq:gaussian-posterior}
\end{equation}
The likelihood uses the raw real-valued network output; it does not
require exact interpolation.  The limit $\sigma\downarrow0$ is taken
with $n,S,N$, and $\lambda$ fixed.  Its limiting optimization problem is
\begin{equation}
 \operatorname*{minimize}_w
 \bigl\{\widehat R_S(z_w)+\lambda Q_N(w)\bigr\}.
 \label{eq:gaussian-objective}
\end{equation}

\subsection{Teacher implementations and a width-independent capacity bound}

\begin{lemma}[Cloning a teacher]
\label{lem:gaussian-cloning}
Suppose that $g$ is a network of the preceding architecture, hidden
width $m$, and depth $H+1$, with every weight of absolute value at most
$B$.  At every width $N\geq m$ it has an exact implementation with
\begin{equation}
 Q_N\leq E:=d(d+1)B^2+2H m^2B^2.
 \label{eq:gaussian-teacher-energy}
\end{equation}
Moreover, $|g(x)|\leq U:=mB$ for every $x$.
\end{lemma}
\begin{proof}
Partition the $N$ neurons in each hidden layer into $m$ groups, of sizes
$k_j\in\{\lfloor N/m\rfloor,\lceil N/m\rceil\}$.  Copy a teacher
first-layer row onto every neuron in its group.  Replace a teacher
internal weight $A^\ell_{ji}$ by $A^\ell_{ji}/k_i$ on every corresponding
edge, and its readout weight $u_j$ by $u_j/k_j$ on each copy.  Every copy
has exactly the teacher activation, so the function is unchanged.
The first-layer energy is at most $d(d+1)B^2$.  Since $k_j/k_i\leq2$,
each internal layer costs at most $2m^2B^2$.  Since
$k_j\geq N/(2m)$, the readout costs at most $2m^2B^2$.  There are
$H-1$ internal layers and one readout.  Finally, $|\tanh|\leq1$ gives
$|g|\leq\sum_j|u_j|\leq mB$.
\end{proof}

Put
\begin{equation}
 D_n=\sqrt{2(d+1)\log(n+1)}+\frac{H-1}{2},\qquad
 \mathcal Z_K=\bigcup_{N\geq1}\{z_w:Q_N(w)\leq K\}.
 \label{eq:gaussian-Dn}
\end{equation}
The union over all learner widths is essential: the following event
controls every width at once.

\begin{lemma}[Capacity of an energy ball]
\label{lem:gaussian-capacity}
For every set of $n$ inputs and every $K>0$, the empirical Rademacher
complexity satisfies
\begin{equation}
 \widehat{\mathfrak R}_S(\mathcal Z_K)
 \leq \left(\frac KH\right)^{H/2}\frac{D_n}{\sqrt n}.
 \label{eq:gaussian-rademacher}
\end{equation}
Consequently, with probability at least $1-\eta$, simultaneously for all
widths and all $w$ with $Q_N(w)\leq K$,
\begin{equation}
 R(\operatorname{clip}z_w)
 \leq\widehat R_S(\operatorname{clip}z_w)
 +8\left(\frac KH\right)^{H/2}\frac{D_n}{\sqrt n}
 +4\sqrt{\frac{\log(1/\eta)}{2n}}.
 \label{eq:gaussian-fixed-energy-deviation}
\end{equation}
\end{lemma}
\begin{proof}
For a bounded set $A\subset\mathbb R^n$, write
$\operatorname{Rad}(A)=n^{-1}\mathbb E_\epsilon
\sup_{a\in A}\sum_i\epsilon_i a_i$.
Its value is unchanged by taking a closed convex hull; it scales
linearly under positive scalar multiplication.  We use the standard
Rademacher contraction inequality, Massart's finite-class bound, and
the Sauer--Shelah bound for halfspaces.

For $t\in[0,1]$, let $\operatorname{clip}_t(u)=\max\{-t,\min\{t,u\}\}$,
and let $\operatorname{Clip}(A)$ contain the vectors
$(\operatorname{clip}_t(a_i))_{i=1}^n$ for $a\in A$ and a single common
$t\in[0,1]$.  The identity
$\operatorname{clip}_t(u)=(|u+t|-|u-t|)/2$, contraction, and symmetry of
the Rademacher signs imply
\begin{equation}
 \operatorname{Rad}(\operatorname{Clip}(A))
 \leq\operatorname{Rad}(A)+\frac1{2\sqrt n}.
 \label{eq:gaussian-variable-clip}
\end{equation}
Indeed, split the supremum between the two absolute-value expressions
and contract each.  The two additional terms sum to
$\mathbb E|\sum_i\epsilon_i|/(2n)\leq1/(2\sqrt n)$.

We also need a positive mixture identity.  For $s>0$ and $|u|\leq1$,
\begin{equation}
 \frac{\tanh(su)}s=
 \int_0^1\operatorname{clip}_t(u)\,
       2s\tanh(st)\operatorname{sech}^2(st)\,dt
 +\operatorname{sech}^2(s)\operatorname{clip}_1(u).
 \label{eq:gaussian-tanh-mixture}
\end{equation}
All weights are nonnegative and have total mass one.  To verify the
identity for $u\geq0$, integrate by parts with
$\psi(t)=\tanh^2(st)$: the right side is
$u-\int_0^u\tanh^2(st)dt=\tanh(su)/s$.  Oddness gives the other half.

On the fixed sample inputs, let $\mathcal F_1$ be the closed convex hull
of the vectors $(\tanh(a\cdot\bar X_i))_{i=1}^n$, with
$a\in\mathbb R^{d+1}$.  Recursively set
$\mathcal F_\ell=\overline{\operatorname{conv}}
\operatorname{Clip}(\mathcal F_{\ell-1})$.  These sets are symmetric,
convex, contain zero, and lie in $[-1,1]^n$.
Define $\rho^1=\boldsymbol1_N$ and
$\rho^\ell=|W^\ell|\rho^{\ell-1}$, with entrywise absolute values.
Induction using \eqref{eq:gaussian-tanh-mixture} shows that the sample
vector of hidden neuron $j$ at layer $\ell$ belongs to
$\rho^\ell_j\mathcal F_\ell$.  A zero value of $\rho^\ell_j$ forces the
neuron to vanish.  For a positive value, divide the preactivation by
$\rho^\ell_j$; it is a convex signed combination of the previous
normalized neurons, hence belongs to $\mathcal F_{\ell-1}$.  The mixture
identity completes the induction.

It follows that the output sample vector belongs to
$P(w)\mathcal F_H$, where
$P(w)=|v|^\top|W^H|\cdots|W^2|\boldsymbol1_N$ (with no matrix factors
when $H=1$).  Cauchy--Schwarz and the Frobenius bound on operator norm give
\[
 P(w)\leq\sqrt N\|v\|_2\prod_{\ell=2}^H\|W^\ell\|_F
 \leq\left(\frac{Q_N(w)}H\right)^{H/2}.
\]
The last inequality is AM--GM over the $H$ energy terms after the first
layer.  The first-layer energy is not needed for this capacity estimate.
By \eqref{eq:gaussian-variable-clip},
$\operatorname{Rad}(\mathcal F_H)\leq
\operatorname{Rad}(\mathcal F_1)+(H-1)/(2\sqrt n)$.

Finally, if a scalar random variable $T$ has distribution function
$(1+\tanh t)/2$, then $\tanh u=\mathbb E\operatorname{sign}(u-T)$.
Thus $\mathcal F_1$ lies in the convex hull of the affine-halfspace
dichotomies of the sample inputs.  There are at most $(n+1)^{d+1}$ such
dichotomies by Sauer--Shelah.  Massart's bound gives
$\operatorname{Rad}(\mathcal F_1)
\leq\sqrt{2(d+1)\log(n+1)/n}$, proving
\eqref{eq:gaussian-rademacher}.

The clipped squared loss is bounded by $4$ and is $4$-Lipschitz in the
raw prediction.  Symmetrization and contraction bound its expected
uniform deviation by eight times
\eqref{eq:gaussian-rademacher}.  Changing one observation changes that
supremum by at most $4/n$, so bounded-difference concentration gives
\eqref{eq:gaussian-fixed-energy-deviation}.
The suprema are measurable: for each width the relevant loss is
continuous in the parameters, and rational parameter vectors in the
open energy ellipsoid are dense in the closed ellipsoid.  Taking the
union over the countably many widths gives a countable determining
class.
\end{proof}

\subsection{Absorbing capacity into the energy penalty}

An estimate based only on the largest possible energy of a minimizer
loses useful information.  In the proof below, the capacity cost of
the selected network is paid by that same network's energy penalty.

\begin{theorem}[Gaussian oracle inequality, at least two hidden layers]
\label{thm:gaussian-absorption}
Suppose $H\geq2$, $0<\delta<1$, and $n\geq1$.  Set
\begin{align}
 a_n&=8(2/H)^{H/2}\frac{D_n}{\sqrt n},&
 \lambda&=a_n^{2/H}
       =\frac2H\left(\frac{8D_n}{\sqrt n}\right)^{2/H},
       \label{eq:gaussian-calibration}\\
 J&=\max\{0,\lceil\log_2(1/\lambda)\rceil\},&
 b_n&=4\sqrt{\frac{\log(2(J+1)/\delta)}{2n}},&
 e_n&=8H^{-H/2}\frac{D_n}{\sqrt n}.
 \label{eq:gaussian-residuals}
\end{align}
Fix a teacher $g$ as in Lemma~\ref{lem:gaussian-cloning}, independently
of the observed sample.  With probability at least $1-\delta$,
simultaneously for every learner width $N\geq m$, every global
minimizer $w$ of \eqref{eq:gaussian-objective} satisfies
\begin{equation}
 R(\operatorname{clip}z_w)
 \leq R(g)+\lambda E+e_n+b_n+h_g,
 \label{eq:gaussian-oracle}
\end{equation}
where one may use either of the following teacher concentration bounds,
chosen in advance (each choice gives its own event of probability at
least $1-\delta$):
\begin{align}
 h_g&=(U+1)^2\sqrt{\frac{\log(2/\delta)}{2n}},
       \label{eq:gaussian-teacher-hoeffding}\\
 h_g&=(U+1)\sqrt{\frac{2R(g)\log(2/\delta)}n}
       +\frac{2(U+1)^2\log(2/\delta)}{3n}.
       \label{eq:gaussian-teacher-bernstein}
\end{align}
On that same event, for each fixed $N\geq m$, the right side of
\eqref{eq:gaussian-oracle} also bounds
\[
 \limsup_{\sigma\downarrow0}
   \mathbb E_{w\sim\Pi_{\sigma,N,S}}R(\operatorname{clip}z_w)
 \quad\hbox{and}\quad
 \limsup_{\sigma\downarrow0}
   R\left(\mathbb E_{w\sim\Pi_{\sigma,N,S}}\operatorname{clip}z_w\right).
\]
The event is common to all widths.  The limit in $\sigma$ is pointwise
in the learner width, with no width-uniform convergence rate asserted.
\end{theorem}
\begin{proof}
The zero network has empirical risk at most one.  Consequently every
global minimizer has energy $q=Q_N(w)\leq1/\lambda$.
Apply Lemma~\ref{lem:gaussian-capacity} at thresholds
$K_j=2^j$, $j=0,\ldots,J$, each with failure probability
$\delta/[2(J+1)]$.  On the resulting event, common to every width,
networks with $q\leq1$ have deviation at most $e_n+b_n$.  For $q>1$,
the least dyadic threshold containing $q$ is less than $2q$, so their
deviation is at most $a_n q^{H/2}+b_n$.  In particular, for every
$q\leq1/\lambda$,
\[
 R(\operatorname{clip}z_w)
 \leq\widehat R_S(\operatorname{clip}z_w)
       +a_nq^{H/2}+e_n+b_n.
\]
Since $H/2\geq1$ and $\lambda^{H/2}=a_n$,
\[
 a_nq^{H/2}\leq a_n\lambda^{1-H/2}q=\lambda q.
\]
Clipping cannot increase squared loss against labels in $[-1,1]$.
Comparison of the minimizing objective with the cloned teacher gives
\[
 R(\operatorname{clip}z_w)
 \leq\widehat R_S(z_w)+\lambda Q_N(w)+e_n+b_n
 \leq\widehat R_S(g)+\lambda E+e_n+b_n.
\]
Hoeffding's inequality gives
\eqref{eq:gaussian-teacher-hoeffding} with failure probability at most
$\delta/2$.  For \eqref{eq:gaussian-teacher-bernstein}, put
$Z=(g(X)-Y)^2$.  Then $0\leq Z\leq(U+1)^2$ and
$\operatorname{Var}(Z)\leq(U+1)^2R(g)$; Bernstein's inequality gives
the stated bound.  Union-bounding with the capacity event proves
\eqref{eq:gaussian-oracle}.

For completeness, at fixed $N,S$ the function
$F(w)=Q_N(w)+\lambda^{-1}\widehat R_S(z_w)$ is continuous and coercive,
so its minimizer set is nonempty and compact.  Write $f_* =\min F$.
For each $\varepsilon>0$, a positive-volume neighborhood of a minimizer
has $F\leq f_*+\varepsilon/2$.  If $\beta=(2\sigma^2)^{-1}$ and
$\beta>\beta_0>0$, then
\[
 \int_{\{F\geq f_*+\varepsilon\}}e^{-\beta F}dw
 \leq e^{-(\beta-\beta_0)(f_*+\varepsilon)}
       \int e^{-\beta_0F}dw.
\]
The final integral is finite because $F\geq Q_N$.  Dividing by the
normalizer lower bound from that neighborhood shows that the posterior
mass of $\{F\geq f_*+\varepsilon\}$ tends to zero.  Coercivity and
continuity then show that the posterior expectation of any bounded
continuous function has upper limit at most its maximum on the
minimizer set.  Apply this to $R(\operatorname{clip}z_w)\in[0,4]$;
its continuity follows by dominated convergence.  Jensen's inequality
gives the assertion for the mean of clipped outputs.
\end{proof}

The leading teacher-dependent term is explicitly
\begin{equation}
 \lambda E=
 \left[\frac{2d(d+1)B^2}{H}+4m^2B^2\right]
 \left(\frac{8D_n}{\sqrt n}\right)^{2/H}.
 \label{eq:gaussian-leading-term}
\end{equation}
Thus, at fixed depth and for a fixed teacher, the excess risk is
$\widetilde O(n^{-1/H})=\widetilde O(n^{-1/(L-1)})$.
For two hidden layers it is $\widetilde O(n^{-1/2})$.
The likelihood variance is $n\lambda\sigma^2$ with the exact
$\lambda$ in \eqref{eq:gaussian-calibration}; it is not a sample-independent
constant.  In particular, up to logarithms this calibration lies at the
power $\lambda\asymp n^{-1/H}$.

\begin{theorem}[One hidden layer]
\label{thm:gaussian-one-hidden}
Suppose $H=1$, and put
\[
 a_n=8\sqrt2\,D_n/\sqrt n,\qquad e_n=8D_n/\sqrt n.
\]
For any $\lambda>0$, define $J,b_n$ by
\eqref{eq:gaussian-residuals}.  With the same assumptions and posterior
interpretation as Theorem~\ref{thm:gaussian-absorption}, the bound is
\begin{equation}
 R(\operatorname{clip}z_w)
 \leq R(g)+\lambda E+\frac{a_n^2}{4\lambda}+e_n+b_n+h_g.
 \label{eq:gaussian-one-hidden-oracle}
\end{equation}
In particular, $\lambda=a_n$ gives an excess of order
$(E+1)D_n/\sqrt n$ plus the concentration terms.  If a deterministic
energy cap $E_0\geq E$ is known, choosing
\begin{equation}
 \lambda=\frac{a_n}{2\sqrt{E_0+1}}
 \label{eq:gaussian-one-hidden-calibration}
\end{equation}
replaces the two energy terms in
\eqref{eq:gaussian-one-hidden-oracle} by at most
$a_n\sqrt{E_0+1}$.
\end{theorem}
\begin{proof}
The same dyadic argument gives deviation at most
$a_n\sqrt q+e_n+b_n$.  After comparing penalized empirical risks, use
\[
 \sup_{q\geq0}\{a_n\sqrt q-\lambda q\}=\frac{a_n^2}{4\lambda}.
\]
The teacher concentration and posterior-limit arguments are unchanged.
For the stated calibration,
$\lambda E+a_n^2/(4\lambda)
=a_n[E/(2\sqrt{E_0+1})+\sqrt{E_0+1}/2]
\leq a_n\sqrt{E_0+1}$.
\end{proof}

\subsection{Comparison with an expanding class of teachers}

\begin{corollary}[Growing teacher width and weight bound]
\label{cor:gaussian-growing-teachers}
Fix $H,d$ and $0<\delta<1$.  Let $\mathcal G_n$ be the class of networks
with $H$ hidden layers, width at most $\lfloor n^\alpha\rfloor$, and
every weight bounded in absolute value by $n^\alpha$.

If $H\geq2$, use \eqref{eq:gaussian-calibration}.  For every
$0\leq\alpha<1/(4H)$, with probability at least $1-\delta$,
simultaneously for all learner widths $N\geq\lfloor n^\alpha\rfloor$,
\begin{equation}
 \limsup_{\sigma\downarrow0}
  \mathbb E_{\Pi_{\sigma,N,S}}R(\operatorname{clip}z_w)
 \leq\min_{g\in\mathcal G_n}R(g)
       +\widetilde O_{H,d,\delta}(n^{-(1/H-4\alpha)}).
 \label{eq:gaussian-growing-deep}
\end{equation}
If $H=1$, use \eqref{eq:gaussian-one-hidden-calibration} with
$E_0=[d(d+1)+2]n^{4\alpha}$.  For every
$0\leq\alpha<1/4$, the corresponding bound is
\begin{equation}
 \limsup_{\sigma\downarrow0}
  \mathbb E_{\Pi_{\sigma,N,S}}R(\operatorname{clip}z_w)
 \leq\min_{g\in\mathcal G_n}R(g)
       +\widetilde O_{d,\delta}(n^{-(1/2-2\alpha)}).
 \label{eq:gaussian-growing-shallow}
\end{equation}
Both assertions also hold for the risk of the posterior mean of clipped
outputs.  The notation $\widetilde O$ suppresses logarithms in $n$;
the preceding theorems give explicit finite-sample constants.
\end{corollary}
\begin{proof}
Narrower teachers embed by zero padding, since $\tanh(0)=0$.
The parameter set defining $\mathcal G_n$ is compact and its population
risk is continuous, so it has a deterministic population-risk minimizer
$g_n$.  The zero network is an available comparator, hence $R(g_n)\leq1$.
Lemma~\ref{lem:gaussian-cloning} gives
\[
 E\leq[d(d+1)+2H]n^{4\alpha},\qquad U\leq n^{2\alpha}.
\]
The Bernstein term \eqref{eq:gaussian-teacher-bernstein} is therefore
$O_\delta(n^{2\alpha-1/2}+n^{4\alpha-1})$.
For $H\geq2$, \eqref{eq:gaussian-leading-term} has order
$\widetilde O_{H,d}(n^{4\alpha-1/H})$, which dominates the preceding
concentration powers.  For $H=1$, the energy-calibrated bound in
Theorem~\ref{thm:gaussian-one-hidden} has order
$\widetilde O_d(n^{2\alpha-1/2})$.  The dyadic concentration terms are
at most logarithmic factors times $n^{-1/2}$ in both cases.
\end{proof}

\paragraph{What the exponent does and does not establish.}
The argument is independent of learner width and allows arbitrary
label noise and arbitrary approximation error.  It still uses a weight
bound, or another bound on the teacher's energy, as well as its width.
It does not give a bound in teacher width alone for arbitrary
real-valued networks.  It also does not make the dependence on depth
polynomial: the leading sample requirement from
\eqref{eq:gaussian-leading-term} raises the teacher-energy and accuracy
factors to the power $H$.

The distinction between clipping before and after posterior averaging
is substantive.  The bounded-loss argument above controls
$\mathbb E_\Pi R(\operatorname{clip}z_w)$ and
$R(\mathbb E_\Pi\operatorname{clip}z_w)$.  By itself it does not control
$R(\mathbb E_\Pi z_w)$ or
$R(\operatorname{clip}(\mathbb E_\Pi z_w))$.
